% Soluciones de la tabla «Binary compounds» del capítulo de formulación (inorganica_apuntes_3_ESO_claude_Libro.tex).
% Se carga desde el fichero de soluciones (\begin{solu} del ejercicio de la tabla) en el libro en español.
% En negro, los datos del enunciado; en rojo, las respuestas.
% Traducción de Soluciones_formulacion_3ESO_ENG.tex (mismos compuestos y números).

\providecommand{\cabeceraSolFormES}[1]{%
    \toprule
    \textbf{Nº} & \textbf{Fórmula} & \textbf{Multiplicadores} & \textbf{Estado de oxidación} \\
    \midrule
    \multicolumn{4}{c}{\textbf{#1}} \\
    \midrule
}

{\small
\begin{center}
    \begin{tabular}{@{}r C{1.9cm} L{5cm} L{5cm}@{}}
        \cabeceraSolFormES{Óxidos}
        1 & \ch{Na2O} & \rojo{Óxido de disodio} & \rojo{Óxido de sodio} \\[2pt]
        2 & \ch{Co2O3} & \rojo{Trióxido de dicobalto} & \rojo{Óxido de cobalto (III)} \\[2pt]
        3 & \ch{CuO} & \rojo{Óxido de cobre (monóxido de cobre)} & \rojo{Óxido de cobre (II)} \\[2pt]
        4 & \ch{Ag2O} & \rojo{Óxido de diplata} & \rojo{Óxido de plata} \\[2pt]
        5 & \ch{CdO} & \rojo{Óxido de cadmio} & \rojo{Óxido de cadmio} \\[2pt]
        6 & \ch{NO} & \rojo{Monóxido de nitrógeno} & \rojo{Óxido de nitrógeno (II)} \\[2pt]
        7 & \ch{CO2} & \rojo{Dióxido de carbono} & \rojo{Óxido de carbono (IV)} \\[2pt]
        8 & \ch{As2O5} & \rojo{Pentaóxido de diarsénico} & \rojo{Óxido de arsénico (V)} \\[2pt]
        9 & \ch{SO3} & \rojo{Trióxido de azufre} & \rojo{Óxido de azufre (VI)} \\[2pt]
        10 & \ch{O5I2} & \rojo{Diyoduro de pentaoxígeno} & -- \\[2pt]
        11 & \ch{O7Br2} & \rojo{Dibromuro de heptaoxígeno} & -- \\[2pt]
        12 & \rojo{\ch{Li2O}} & Óxido de dilitio & \rojo{Óxido de litio} \\[2pt]
        13 & \rojo{\ch{MgO}} & \rojo{Óxido de magnesio} & Óxido de magnesio \\[2pt]
        14 & \rojo{\ch{K2O}} & \rojo{Óxido de dipotasio} & Óxido de potasio \\[2pt]
        15 & \rojo{\ch{SrO}} & \rojo{Óxido de estroncio} & Óxido de estroncio \\[2pt]
        16 & \rojo{\ch{FeO}} & \rojo{Óxido de hierro (monóxido de hierro)} & Óxido de hierro (II) \\[2pt]
        17 & \rojo{\ch{PbO2}} & Dióxido de plomo & \rojo{Óxido de plomo (IV)} \\[2pt]
        18 & \rojo{\ch{Au2O3}} & \rojo{Trióxido de dioro} & Óxido de oro (III) \\[2pt]
        19 & \rojo{\ch{Al2O3}} & Trióxido de dialuminio & \rojo{Óxido de aluminio} \\[2pt]
        20 & \rojo{\ch{SnO}} & \rojo{Óxido de estaño (monóxido de estaño)} & Óxido de estaño (II) \\[2pt]
        21 & \rojo{\ch{ZnO}} & \rojo{Óxido de cinc} & Óxido de cinc \\[2pt]
        22 & \rojo{\ch{HgO}} & \rojo{Óxido de mercurio (monóxido de mercurio)} & Óxido de mercurio (II) \\[2pt]
        23 & \rojo{\ch{N2O4}} & Tetraóxido de dinitrógeno & \rojo{Óxido de nitrógeno (IV)} \\[2pt]
        24 & \rojo{\ch{Co2O3}} & \rojo{Trióxido de dicobalto} & Óxido de cobalto (III) \\[2pt]
        25 & \rojo{\ch{N2O5}} & Pentaóxido de dinitrógeno & \rojo{Óxido de nitrógeno (V)} \\[2pt]
        26 & \rojo{\ch{BaO}} & \rojo{Óxido de bario} & Óxido de bario \\[2pt]
        27 & \rojo{\ch{SO3}} & Trióxido de azufre & \rojo{Óxido de azufre (VI)} \\[2pt]
        28 & \rojo{\ch{Cr2O3}} & Trióxido de dicromo & \rojo{Óxido de cromo (III)} \\[2pt]
        29 & \rojo{\ch{MnO2}} & \rojo{Dióxido de manganeso} & Óxido de manganeso (IV) \\[2pt]
        30 & \rojo{\ch{N2O}} & Monóxido de dinitrógeno & \rojo{Óxido de nitrógeno (I)} \\[2pt]
        31 & \rojo{\ch{PtO2}} & \rojo{Dióxido de platino} & Óxido de platino (IV) \\[2pt]
        32 & \rojo{\ch{O3Br2}} & Dibromuro de trioxígeno & -- \\
        \bottomrule
    \end{tabular}
\end{center}

\noindent En 10, 11 y 32 el halógeno va a la derecha (es más electronegativo que el oxígeno), así que solo se nombran con multiplicadores.

\begin{center}
    \begin{tabular}{@{}r C{1.9cm} L{5cm} L{5cm}@{}}
        \cabeceraSolFormES{Hidruros}
        33 & \ch{NaH} & \rojo{Hidruro de sodio} & \rojo{Hidruro de sodio} \\[2pt]
        34 & \ch{CaH2} & \rojo{Dihidruro de calcio} & \rojo{Hidruro de calcio} \\[2pt]
        35 & \ch{CoH2} & \rojo{Dihidruro de cobalto} & \rojo{Hidruro de cobalto (II)} \\[2pt]
        36 & \ch{CuH2} & \rojo{Dihidruro de cobre} & \rojo{Hidruro de cobre (II)} \\[2pt]
        37 & \ch{SnH4} & \rojo{Tetrahidruro de estaño} & \rojo{Hidruro de estaño (IV)} \\[2pt]
        38 & \ch{LiH} & \rojo{Hidruro de litio} & \rojo{Hidruro de litio} \\[2pt]
        39 & \ch{HgH2} & \rojo{Dihidruro de mercurio} & \rojo{Hidruro de mercurio (II)} \\[2pt]
        40 & \ch{NiH3} & \rojo{Trihidruro de níquel} & \rojo{Hidruro de níquel (III)} \\[2pt]
        41 & \ch{AlH3} & \rojo{Trihidruro de aluminio} & \rojo{Hidruro de aluminio} \\[2pt]
        42 & \ch{H4Pb} & \rojo{Tetrahidruro de plomo} & \rojo{Hidruro de plomo (IV)} \\[2pt]
        43 & \ch{H3Mn} & \rojo{Trihidruro de manganeso} & \rojo{Hidruro de manganeso (III)} \\[2pt]
        44 & \rojo{\ch{MgH2}} & \rojo{Dihidruro de magnesio} & Hidruro de magnesio \\[2pt]
        45 & \rojo{\ch{NaH}} & \rojo{Hidruro de sodio} & Hidruro de sodio \\[2pt]
        46 & \rojo{\ch{SrH2}} & \rojo{Dihidruro de estroncio} & Hidruro de estroncio \\[2pt]
        47 & \rojo{\ch{CrH2}} & \rojo{Dihidruro de cromo} & Hidruro de cromo (II) \\[2pt]
        48 & \rojo{\ch{FeH3}} & Trihidruro de hierro & \rojo{Hidruro de hierro (III)} \\[2pt]
        49 & \rojo{\ch{PbH4}} & \rojo{Tetrahidruro de plomo} & Hidruro de plomo (IV) \\[2pt]
        50 & \rojo{\ch{ZnH2}} & Dihidruro de cinc & \rojo{Hidruro de cinc} \\[2pt]
        51 & \rojo{\ch{CdH2}} & \rojo{Dihidruro de cadmio} & Hidruro de cadmio \\[2pt]
        52 & \rojo{\ch{CuH}} & \rojo{Hidruro de cobre} & Hidruro de cobre (I) \\[2pt]
        53 & \rojo{\ch{AuH3}} & \rojo{Trihidruro de oro} & Hidruro de oro (III) \\[2pt]
        54 & \rojo{\ch{SnH2}} & \rojo{Dihidruro de estaño} & Hidruro de estaño (II) \\[2pt]
        55 & \rojo{\ch{KH}} & Hidruro de potasio & \rojo{Hidruro de potasio} \\[2pt]
        56 & \rojo{\ch{CoH2}} & \rojo{Dihidruro de cobalto} & Hidruro de cobalto (II) \\[2pt]
        57 & \rojo{\ch{PtH4}} & \rojo{Tetrahidruro de platino} & Hidruro de platino (IV) \\[2pt]
        58 & \rojo{\ch{BaH2}} & Dihidruro de bario & \rojo{Hidruro de bario} \\[2pt]
        59 & \rojo{\ch{AgH}} & Hidruro de plata & \rojo{Hidruro de plata} \\[2pt]
        60 & \rojo{\ch{AuH}} & \rojo{Hidruro de oro} & Hidruro de oro (I) \\
        \bottomrule
    \end{tabular}
\end{center}

\noindent En 42 y 43 lo habitual es escribir primero el metal: \ch{PbH4} y \ch{MnH3}.

\begin{center}
    \begin{tabular}{@{}r C{1.9cm} L{5cm} L{5cm}@{}}
        \cabeceraSolFormES{Hidruros no metálicos (hidrógeno con un no metal de los grupos 16 y 17)}
        61 & \ch{HCl} & \rojo{Cloruro de hidrógeno} & \rojo{Cloruro de hidrógeno} \\[2pt]
        62 & \ch{HBr} & \rojo{Bromuro de hidrógeno} & \rojo{Bromuro de hidrógeno} \\[2pt]
        63 & \ch{HI} & \rojo{Yoduro de hidrógeno} & \rojo{Yoduro de hidrógeno} \\[2pt]
        64 & \ch{HF} & \rojo{Fluoruro de hidrógeno} & \rojo{Fluoruro de hidrógeno} \\[2pt]
        65 & \ch{H2S} & \rojo{Sulfuro de dihidrógeno} & \rojo{Sulfuro de hidrógeno} \\[2pt]
        66 & \ch{H2Te} & \rojo{Telururo de dihidrógeno} & \rojo{Telururo de hidrógeno} \\[2pt]
        67 & \ch{H2Se} & \rojo{Seleniuro de dihidrógeno} & \rojo{Seleniuro de hidrógeno} \\[2pt]
        68 & \rojo{\ch{HCl}} & Cloruro de hidrógeno & \rojo{Cloruro de hidrógeno} \\[2pt]
        69 & \rojo{\ch{HBr}} & \rojo{Bromuro de hidrógeno} & Bromuro de hidrógeno \\[2pt]
        70 & \rojo{\ch{HI}} & \rojo{Yoduro de hidrógeno} & Yoduro de hidrógeno \\[2pt]
        71 & \rojo{\ch{HF}} & \rojo{Fluoruro de hidrógeno} & Fluoruro de hidrógeno \\[2pt]
        72 & \rojo{\ch{H2S}} & Sulfuro de dihidrógeno & \rojo{Sulfuro de hidrógeno} \\[2pt]
        73 & \rojo{\ch{H2Te}} & \rojo{Telururo de dihidrógeno} & Telururo de hidrógeno \\[2pt]
        74 & \rojo{\ch{H2Se}} & \rojo{Seleniuro de dihidrógeno} & Seleniuro de hidrógeno \\
        \bottomrule
    \end{tabular}
\end{center}

\noindent El hidrógeno solo tiene un estado de oxidación (+1), así que el segundo nombre es el mismo sin prefijos.

\begin{center}
    \begin{tabular}{@{}r C{1.9cm} L{5cm} L{5cm}@{}}
        \cabeceraSolFormES{Sales binarias}
        75 & \ch{MgCl2} & \rojo{Dicloruro de magnesio} & \rojo{Cloruro de magnesio} \\[2pt]
        76 & \ch{BaBr2} & \rojo{Dibromuro de bario} & \rojo{Bromuro de bario} \\[2pt]
        77 & \ch{Na3N} & \rojo{Nitruro de trisodio} & \rojo{Nitruro de sodio} \\[2pt]
        78 & \ch{K2S} & \rojo{Sulfuro de dipotasio} & \rojo{Sulfuro de potasio} \\[2pt]
        79 & \ch{CaI2} & \rojo{Diyoduro de calcio} & \rojo{Yoduro de calcio} \\[2pt]
        80 & \ch{FeBr3} & \rojo{Tribromuro de hierro} & \rojo{Bromuro de hierro (III)} \\[2pt]
        81 & \ch{PbS2} & \rojo{Disulfuro de plomo} & \rojo{Sulfuro de plomo (IV)} \\[2pt]
        82 & \ch{Au2S3} & \rojo{Trisulfuro de dioro} & \rojo{Sulfuro de oro (III)} \\[2pt]
        83 & \ch{HgF2} & \rojo{Difluoruro de mercurio} & \rojo{Fluoruro de mercurio (II)} \\[2pt]
        84 & \ch{AlCl3} & \rojo{Tricloruro de aluminio} & \rojo{Cloruro de aluminio} \\[2pt]
        85 & \ch{SnBr2} & \rojo{Dibromuro de estaño} & \rojo{Bromuro de estaño (II)} \\[2pt]
        86 & \ch{CuI} & \rojo{Yoduro de cobre} & \rojo{Yoduro de cobre (I)} \\[2pt]
        87 & \ch{CI4} & \rojo{Tetrayoduro de carbono} & \rojo{Yoduro de carbono (IV)} \\[2pt]
        88 & \ch{NF3} & \rojo{Trifluoruro de nitrógeno} & \rojo{Fluoruro de nitrógeno (III)} \\[2pt]
        89 & \ch{IBr} & \rojo{Bromuro de yodo} & \rojo{Bromuro de yodo (I)} \\[2pt]
        90 & \ch{SeBr4} & \rojo{Tetrabromuro de selenio} & \rojo{Bromuro de selenio (IV)} \\[2pt]
        91 & \ch{SeCl6} & \rojo{Hexacloruro de selenio} & \rojo{Cloruro de selenio (VI)} \\[2pt]
        92 & \ch{AsF5} & \rojo{Pentafluoruro de arsénico} & \rojo{Fluoruro de arsénico (V)} \\[2pt]
        93 & \ch{CCl4} & \rojo{Tetracloruro de carbono} & \rojo{Cloruro de carbono (IV)} \\[2pt]
        94 & \ch{SiF4} & \rojo{Tetrafluoruro de silicio} & \rojo{Fluoruro de silicio} \\
        \bottomrule
    \end{tabular}
\end{center}

\begin{center}
    \begin{tabular}{@{}r C{1.9cm} L{5cm} L{5cm}@{}}
        \cabeceraSolFormES{Sales binarias (continuación)}
        95 & \rojo{\ch{LiCl}} & \rojo{Cloruro de litio} & Cloruro de litio \\[2pt]
        96 & \rojo{\ch{MgI2}} & Diyoduro de magnesio & \rojo{Yoduro de magnesio} \\[2pt]
        97 & \rojo{\ch{BeBr2}} & \rojo{Dibromuro de berilio} & Bromuro de berilio \\[2pt]
        98 & \rojo{\ch{K2S}} & Sulfuro de dipotasio & \rojo{Sulfuro de potasio} \\[2pt]
        99 & \rojo{\ch{FeF2}} & \rojo{Difluoruro de hierro} & Fluoruro de hierro (II) \\[2pt]
        100 & \rojo{\ch{PbCl4}} & Tetracloruro de plomo & \rojo{Cloruro de plomo (IV)} \\[2pt]
        101 & \rojo{\ch{Au2S3}} & \rojo{Trisulfuro de dioro} & Sulfuro de oro (III) \\[2pt]
        102 & \rojo{\ch{SnS}} & \rojo{Sulfuro de estaño (monosulfuro de estaño)} & Sulfuro de estaño (II) \\[2pt]
        103 & \rojo{\ch{ZnBr2}} & Dibromuro de cinc & \rojo{Bromuro de cinc} \\[2pt]
        104 & \rojo{\ch{HgTe}} & \rojo{Telururo de mercurio} & Telururo de mercurio (II) \\[2pt]
        105 & \rojo{\ch{CoI2}} & \rojo{Diyoduro de cobalto} & Yoduro de cobalto (II) \\[2pt]
        106 & \rojo{\ch{CaCl2}} & \rojo{Dicloruro de calcio} & Cloruro de calcio \\[2pt]
        107 & \rojo{\ch{Al2S3}} & Trisulfuro de dialuminio & \rojo{Sulfuro de aluminio} \\[2pt]
        108 & \rojo{\ch{Cr2S3}} & Trisulfuro de dicromo & \rojo{Sulfuro de cromo (III)} \\[2pt]
        109 & \rojo{\ch{MnBr2}} & \rojo{Dibromuro de manganeso} & Bromuro de manganeso (II) \\[2pt]
        110 & \rojo{\ch{Na2Te}} & Telururo de disodio & \rojo{Telururo de sodio} \\[2pt]
        111 & \rojo{\ch{ZnCl2}} & \rojo{Dicloruro de cinc} & Cloruro de cinc \\[2pt]
        112 & \rojo{\ch{Ag2S}} & Sulfuro de diplata & \rojo{Sulfuro de plata} \\[2pt]
        113 & \rojo{\ch{BrF3}} & Trifluoruro de bromo & \rojo{Fluoruro de bromo (III)} \\[2pt]
        114 & \rojo{\ch{BrCl}} & \rojo{Cloruro de bromo} & Cloruro de bromo (I) \\[2pt]
        115 & \rojo{\ch{SF6}} & Hexafluoruro de azufre & \rojo{Fluoruro de azufre (VI)} \\[2pt]
        116 & \rojo{\ch{CS2}} & \rojo{Disulfuro de carbono} & Sulfuro de carbono (IV) \\
        \bottomrule
    \end{tabular}
\end{center}

\begin{center}
    \begin{tabular}{@{}r C{1.9cm} L{5cm} L{5cm}@{}}
        \cabeceraSolFormES{Hidróxidos}
        117 & \ch{CsOH} & \rojo{Hidróxido de cesio} & \rojo{Hidróxido de cesio} \\[2pt]
        118 & \ch{KOH} & \rojo{Hidróxido de potasio} & \rojo{Hidróxido de potasio} \\[2pt]
        119 & \ch{Be(OH)2} & \rojo{Dihidróxido de berilio} & \rojo{Hidróxido de berilio} \\[2pt]
        120 & \ch{Fe(OH)3} & \rojo{Trihidróxido de hierro} & \rojo{Hidróxido de hierro (III)} \\[2pt]
        121 & \ch{AgOH} & \rojo{Hidróxido de plata} & \rojo{Hidróxido de plata} \\[2pt]
        122 & \ch{Al(OH)3} & \rojo{Trihidróxido de aluminio} & \rojo{Hidróxido de aluminio} \\[2pt]
        123 & \ch{NH4OH} & \rojo{Hidróxido de amonio} & \rojo{Hidróxido de amonio} \\[2pt]
        124 & \ch{Cd(OH)2} & \rojo{Dihidróxido de cadmio} & \rojo{Hidróxido de cadmio} \\[2pt]
        125 & \ch{Au(OH)3} & \rojo{Trihidróxido de oro} & \rojo{Hidróxido de oro (III)} \\[2pt]
        126 & \ch{Mn(OH)3} & \rojo{Trihidróxido de manganeso} & \rojo{Hidróxido de manganeso (III)} \\[2pt]
        127 & \ch{Pt(OH)2} & \rojo{Dihidróxido de platino} & \rojo{Hidróxido de platino (II)} \\[2pt]
        128 & \rojo{\ch{Ni(OH)2}} & Dihidróxido de níquel & \rojo{Hidróxido de níquel (II)} \\[2pt]
        129 & \rojo{\ch{Ni(OH)3}} & \rojo{Trihidróxido de níquel} & Hidróxido de níquel (III) \\[2pt]
        130 & \rojo{\ch{NaOH}} & \rojo{Hidróxido de sodio} & Hidróxido de sodio \\[2pt]
        131 & \rojo{\ch{Ca(OH)2}} & \rojo{Dihidróxido de calcio} & Hidróxido de calcio \\[2pt]
        132 & \rojo{\ch{Mg(OH)2}} & Dihidróxido de magnesio & \rojo{Hidróxido de magnesio} \\[2pt]
        133 & \rojo{\ch{Hg(OH)2}} & \rojo{Dihidróxido de mercurio} & Hidróxido de mercurio (II) \\[2pt]
        134 & \rojo{\ch{Fe(OH)2}} & Dihidróxido de hierro & \rojo{Hidróxido de hierro (II)} \\[2pt]
        135 & \rojo{\ch{Zn(OH)2}} & \rojo{Dihidróxido de cinc} & Hidróxido de cinc \\[2pt]
        136 & \rojo{\ch{Cu(OH)2}} & Dihidróxido de cobre & \rojo{Hidróxido de cobre (II)} \\[2pt]
        137 & \rojo{\ch{Pb(OH)4}} & \rojo{Tetrahidróxido de plomo} & Hidróxido de plomo (IV) \\[2pt]
        138 & \rojo{\ch{Sn(OH)2}} & \rojo{Dihidróxido de estaño} & Hidróxido de estaño (II) \\[2pt]
        139 & \rojo{\ch{Ni(OH)3}} & \rojo{Trihidróxido de níquel} & Hidróxido de níquel (III) \\[2pt]
        140 & \rojo{\ch{Cr(OH)3}} & Trihidróxido de cromo & \rojo{Hidróxido de cromo (III)} \\[2pt]
        141 & \rojo{\ch{Cr(OH)2}} & \rojo{Dihidróxido de cromo} & Hidróxido de cromo (II) \\
        \bottomrule
    \end{tabular}
\end{center}

\begin{center}
    \begin{tabular}{@{}r C{1.9cm} L{5cm} L{5cm}@{}}
        \cabeceraSolFormES{Peróxidos}
        142 & \ch{Na2O2} & \rojo{Peróxido de sodio} & -- \\[2pt]
        143 & \ch{K2O2} & \rojo{Peróxido de potasio} & -- \\[2pt]
        144 & \ch{SrO2} & \rojo{Peróxido de estroncio} & -- \\[2pt]
        145 & \ch{CaO2} & \rojo{Peróxido de calcio} & -- \\[2pt]
        146 & \ch{Rb2O2} & \rojo{Peróxido de rubidio} & -- \\[2pt]
        147 & \ch{CuO2} & \rojo{Peróxido de cobre (II)} & -- \\[2pt]
        148 & \ch{MgO2} & \rojo{Peróxido de magnesio} & -- \\[2pt]
        149 & \ch{BaO2} & \rojo{Peróxido de bario} & -- \\
        \midrule
        \multicolumn{4}{c}{\textbf{Hidruros de los grupos 16 y 17 en disolución acuosa}} \\
        \midrule
        150 & \ch{HCl(ac)} & \rojo{Ácido clorhídrico} & -- \\[2pt]
        151 & \ch{HBr(ac)} & \rojo{Ácido bromhídrico} & -- \\[2pt]
        152 & \ch{HI(ac)} & \rojo{Ácido yodhídrico} & -- \\[2pt]
        153 & \ch{HF(ac)} & \rojo{Ácido fluorhídrico} & -- \\[2pt]
        154 & \ch{H2S(ac)} & \rojo{Ácido sulfhídrico} & -- \\[2pt]
        155 & \ch{H2Te(ac)} & \rojo{Ácido telurhídrico} & -- \\[2pt]
        156 & \ch{H2Se(ac)} & \rojo{Ácido selenhídrico} & -- \\
        \bottomrule
    \end{tabular}
\end{center}
}
